\begingroup\footnotesize\raggedright
\begin{itemize}[nosep,leftmargin=*]
\item \prop{infl}: \emph{Structural reasoning.} \bprint{INFL(PrtDriver, Problem1|PTROFFLINE)∧ P(Problem1=No\_Output|PrtDriver=No) ≥0.5} The printer driver influences the no-output symptom and the conditional failure probability is at least 0.5.
\item \prop{map\_conf}: \emph{MAP and MPE.} \bprint{P(MAP(NetPrint, NetOK|Problem1=No\_Output)|Problem1=No\_Output) >=0.5-0.01} The result is the most probable network print explanation accounts for at least 0.49 of the posterior mass.
\item \prop{map\_update}: \emph{MAP and MPE.} \bprint{P(MAP(NetPrint, NetOK|Problem1=No\_Output) [AppOK=Incorrect\_Corrupt↦0.5]↔MAP(NetPrint, NetOK|Problem1=No\_Output))=1} The most probable network print setting is unchanged with probability 1: no local printer and network OK.
\item \prop{prob\_delta}: \emph{What-if analysis.} \bprint{P(Problem1=No\_Output) [PrtCbl=Loose↦0.5]− P(Problem1=No\_Output)} The result is an increase of the failure probability by 0.1680.
\item \prop{prob\_offline}: \emph{What-if analysis.} \bprint{P(Problem1=No\_Output|PTROFFLINE=Offline∨ ¬(PrtOn=Yes)) [PrtCbl=Loose↦0.5] [PrtMem=Less\_than\_2Mb↦0.4] ≥0.6} The probability of no output is at least 0.6 under both parameter updates.
\end{itemize}
\endgroup
